% FAC-EN-U0076
% Source: fac_body.tex lines 3149--3196.
% Authority: PDF physical pp. 78--79; printed pp. 273--274
% (Chapter III no. 78 complete; hard stop before section 6/no. 79).
% French control: FAC-FR-C0041 corrects the printed basis-degree index.

\medskip
\phantomsection\label{fac:III:5:no:78}
\textbf{78. Complete intersections.}

\phantomsection\label{fac:III:5:no:78:def:complete-intersection}
A subvariety $V$ of dimension $p$ in the projective space
$\mathbf P_r(K)$ is a \textit{complete intersection} if the ideal
$I(V)$ of polynomials vanishing on $V$ has a system of $r-p$ generators
\phantomsection\label{fac:III:5:no:78:external-citation:Macaulay-theorem}
$P_1,\ldots,P_{r-p}$; in that case, all the irreducible components of
$V$ have dimension $p$, by Macaulay's theorem (cf. [9], no.~17). It is
well known that such a variety is $p$-fold of the first kind, which
already implies that $H^q(V,\mathcal O_V(n))=0$ for $0<q<p$, as we have
\hyperref[fac:III:5:no:77:prop:3:ii]{just seen}. We shall determine $H^p(V,\mathcal O_V(n))$ in terms of the
degrees $m_1,\ldots,m_{r-p}$ of the homogeneous polynomials
$P_1,\ldots,P_{r-p}$.

Let $S(V)=S/I(V)$ be the projective coordinate ring of $V$. By
\hyperref[fac:III:4:no:72:thm:1]{Theorem~1 of no.~72}, everything
reduces to determining the $S$-module
$\operatorname{Ext}_S^{r-p}(S(V),\Omega)$. There is a resolution
analogous to that of
\hyperref[fac:III:5:no:75:formula:resolution]{no.~75}: take for $L^q$
the free graded $S$-module having as basis the elements
$e\langle i_1\cdots i_q\rangle$ corresponding to the sequences
$(i_1,\ldots,i_q)$ such that
$1\leq i_1<i_2<\cdots<i_q\leq r-p$, and having degrees
\phantomsection\label{fac:III:5:no:78:def:Lq}
$\sum_{j=1}^{j=q}m_{i_j}$;\footnote{The printed French has
$\sum_{j=1}^{j=q}m_j$. The basis element depends on the selected indices
$i_1,\ldots,i_q$, the differential below uses $P_{i_j}$, and the next
paragraph explicitly gives degree
$m_{i_1}+\cdots+m_{i_q}+n-r-1$. Corrected French and English therefore
restore the missing $i_j$ in the subscript; diplomatic French preserves
the print.} for $L^0$, take $S$. Set
\phantomsection\label{fac:III:5:no:78:formula:differential}
\[
d(e\langle i_1\cdots i_q\rangle)
=\sum_{j=1}^{j=q}(-1)^jP_{i_j}\cdot
e\langle i_1\cdots\widehat{\imath}_j\cdots i_q\rangle,
\qquad
d(e\langle i\rangle)=P_i.
\]

The sequence
\phantomsection\label{fac:III:5:no:78:formula:resolution}
\[
0\longrightarrow L^{r-p}\overset d\longrightarrow\cdots
\overset d\longrightarrow L^0\longrightarrow S(V)\longrightarrow0
\]
is exact ([6], Chapter~VIII, Proposition~4.3). It follows that the
$\operatorname{Ext}_S^q(S(V),\Omega)$ are the cohomology groups of the
complex formed by the $\operatorname{Hom}_S(L^q,\Omega)$. An element of
$\operatorname{Hom}_S(L^q,\Omega)_n$ may be identified with a system
$f\langle i_1\cdots i_q\rangle$, where the
$f\langle i_1\cdots i_q\rangle$ are homogeneous polynomials of degrees
$m_{i_1}+\cdots+m_{i_q}+n-r-1$. Under this identification, the
coboundary operator is given by the usual formula
\phantomsection\label{fac:III:5:no:78:formula:coboundary}
\[
(df)\langle i_1\cdots i_{q+1}\rangle
=\sum_{j=1}^{j=q+1}(-1)^jP_{i_j}\cdot
f\langle i_1\cdots\widehat{\imath}_j\cdots i_{q+1}\rangle.
\]

Macaulay's theorem \hyperref[fac:III:5:no:78:external-citation:Macaulay-theorem]{cited above} shows that the conditions of [11] are
satisfied, and one recovers the fact that
\phantomsection\label{fac:III:5:no:78:formula:Ext-concentration}
$\operatorname{Ext}_S^q(S(V),\Omega)=0$ for $q\neq r-p$. On the other
hand, $\operatorname{Ext}_S^{r-p}(S(V),\Omega)_n$ is isomorphic to the
subspace of $S(V)$ formed by the homogeneous elements of degree $N+n$,
where
\phantomsection\label{fac:III:5:no:78:def:N}
\[
N=\sum_{i=1}^{i=r-p}m_i-r-1.
\]
Taking \hyperref[fac:III:4:no:72:thm:1]{Theorem~1 of no.~72} into
account, one obtains:

\phantomsection\label{fac:III:5:no:78:prop:5}
\textsc{Proposition 5.} \textit{Let $V$ be a complete intersection
defined by homogeneous polynomials $P_1,\ldots,P_{r-p}$ of degrees
$m_1,\ldots,m_{r-p}$.}

\phantomsection\label{fac:III:5:no:78:prop:5:a}
(a) \textit{The map
$\alpha:S(V)_n\to H^0(V,\mathcal O_V(n))$ is bijective for every
$n\in\mathbf Z$.}

\phantomsection\label{fac:III:5:no:78:prop:5:b}
(b) \textit{$H^q(V,\mathcal O_V(n))=0$ for $0<q<p$ and every
$n\in\mathbf Z$.}

\phantomsection\label{fac:III:5:no:78:prop:5:c}
(c) \textit{$H^p(V,\mathcal O_V(n))$ is isomorphic to the dual vector
space of $H^0(V,\mathcal O_V(N-n))$, where
$N=\sum_{i=1}^{i=r-p}m_i-r-1$.}

\phantomsection\label{fac:III:5:no:78:remark}
In particular, one should note that $H^p(V,\mathcal O_V)$ can vanish
only if $N<0$.
